% ch00 译文：原文第 1 节 1a（main.tex 第 171–214 行）
% [原文 p.2]
\section{引言}

镜像对称现象最初由物理学家发现，后来以多种不同形式被数学地表述。它先被理解为两个 Calabi--Yau 流形具有互为镜像的霍奇菱形，随后又被理解为周期积分与 Gromov--Witten 不变量之间的关系。1994 年，Kontsevich \cite{Kontsevich95} 提出了一个系统理解镜像对称现象的框架：把它看成三角范畴之间的等价，其中一侧来自代数簇上层的凝聚复形，另一侧来自辛流形的 Lagrangian 子流形。这就是所谓的 \emph{同调镜像对称} 纲领。用现代语言更精确地说，给定一对期望互为镜像的辛流形 $(X,\omega)$ 与代数簇 $(Y,J)$，我们期望 $A_{\infty}$-范畴之间存在拟等价
\[
    \mathcal{F} (X,\omega)\simeq\Coh (Y,J)
\]
它连接辛流形 $(X,\omega)$ 的某一版本 \emph{Fukaya 范畴} 与 $Y$ 上凝聚层复形的 dg 范畴。一旦知道同调镜像对称，就应当能通过纯代数构造 \cite{KKP08} 导出镜像对称的计数版本。

\begin{conceptbox}{同调镜像对称（HMS）}
Kontsevich 1994 年的猜想：一对镜像的辛流形 $X$ 与代数簇 $Y$ 之间，存在 $A_{\infty}$-范畴的拟等价 $\mathcal{F}(X)\simeq\Coh(Y)$。左侧（A 模型）只依赖 $X$ 的辛结构——对象是 Lagrangian 子流形连同局部系统，态射由 Floer 上同调给出；右侧（B 模型）只依赖 $Y$ 的复结构——对象是凝聚层复形，态射是导出 Hom。两侧分别来自 IIA 型与 IIB 型弦论的 D-膜。猜想的威力在于把「几何上互为镜像」这一模糊说法翻译成精确的范畴等价；已知情形包括椭圆曲线、四次曲面、环簇 Fano 流形与对数 Calabi--Yau 曲面等。
\par\smallskip\noindent{\footnotesize 参考：nLab，\emph{homological mirror symmetry}；Kontsevich, \emph{Homological algebra of mirror symmetry}, 1994。}
\end{conceptbox}

% [原文 p.3]
然而，同调镜像对称纲领并未回答镜像对如何构造的问题；相应的几何构造由 Strominger--Yau--Zaslow \cite{SYZ96} 于 1996 年提出。粗略地说，设辛流形 $(X,\omega)$ 具有纤维为 Lagrangian 环面的纤维化 $X\to B$，则可通过取纤维环面的对偶构造镜像空间 $\check{X}\to B$。在好的情形下，对偶空间上有自然的复结构 $J$，于是得到候选的镜像空间 $(\check{X},J)$。这一构造随后被 Gross 与 Siebert 的著名工作 \cite{GS06,GS10,GS26} 推广到 $X\to B$ 为可能带椭圆奇异纤维的 Lagrangian 环面纤维化的情形。

\begin{yizhu}[SYZ 对偶究竟对偶了什么]
纤维环面 $T^n=\mathbb{R}^n/\mathbb{Z}^n$ 上带有由 $\omega$ 与 Lagrangian 纤维化给出的平坦坐标；对偶环面 $\check{T}^n=\operatorname{Hom}(T^n,U(1))$ 继承对偶格，而复化辛形式在对偶侧给出的平坦坐标恰好定义一个复结构。所以「对偶」改变的是格：A 侧（辛体积）的数据变成 B 侧（复模）的参数。纤维化的奇异纤维处不能再这样逐纤维对偶，需换成仿射流形与辛奇点，这正是 Gross--Siebert 纲领处理的内容。
\end{yizhu}

在 \cite{Fukaya05} 中，Fukaya 提出了一条途径：当 $\check{X}$ 由 $X$ 经类 SYZ 构造得到时，构造镜像函子 $\mathcal{F} (X,\omega)\to\Coh (\check{X},J)$。粗略地说，我们并不天真地取对偶的 Lagrangian 纤维化，而是考察 Fukaya 范畴 $\mathcal{F} (X,\omega)$ 中 Lagrangian 环面（及其「闭包」）对象的模空间，于是镜像函子的存在几乎就是定义。沿着这一纲领，Abouzaid \cite{Abouzaid21} 在 $X\to B$ 为光滑环面纤维化且 $\pi_2(B)=0$ 时取得了重要的部分进展；最近 Bai--Haney \cite{BH26} 借助内蕴镜像对称，在某些额外技术假设下验证了 $X\to B$ 属于 Gross--Siebert 纲领的情形。

本文考察 $A$ 模型与 $B$ 模型互换的镜像对 $(X,\check{X})$。换言之，我们希望理解下述问题：

\begin{ques}\label{Q:HMS-toric-degeneration}
    设 $(X_t)_{t\in\mathbb{C}}$ 为环簇退化族，$(Y,\omega_{p(t)})$ 为 $(X_t, J_t)$ 的镜像辛流形。$A_{\infty}$-范畴之间是否存在拟等价
\[
        \Coh (X_t,J_t)\simeq\mathcal{F} (Y,\omega_{p(t)})?
    \]
\end{ques}

\begin{yizhu}[为什么把 A、B 两侧互换]
通常的同调镜像对称陈述写作 $\mathcal{F}(X)\simeq\Coh(Y)$，其中辛侧是 A 模型、复几何侧是 B 模型。问题 \ref{Q:HMS-toric-degeneration} 只是把同一对镜像换一个立足点：退化族 $(X_t)$ 是复几何侧（B 模型），它的镜像 $(Y,\omega)$ 才是辛侧（A 模型），所以等价写成 $\Coh(X_t)\simeq\mathcal{F}(Y)$。这样写的好处是：$t$ 变化时变的是 $X_t$ 的复结构，而 $\omega_{p(t)}$ 随之变化——退化、紧化这些几何问题因此被搬到辛侧。
\end{yizhu}

注意，当 $X_t\equiv X_{\Sigma}$ 为矩多面体 $\Sigma$ 对应的环簇时，问题 \ref{Q:HMS-toric-degeneration} 已由 Fang--Liu--Treumann--Zaslow \cite{FLTZ12} 沿 Bondal \cite{Bondal06} 的早期想法以及 Bondal--Ruan 的未发表工作给出解答：此时镜像 $(Y,\omega)=\left((\mathbb{C}^{\ast})^n,\Lambda\right)$ 是类 Landau--Ginzburg 模型的辛流形，更严格地说是 \emph{Liouville 扇形}，其中所谈的 Fukaya 范畴是 \emph{部分包裹 Fukaya 范畴}，所用语言见 \cite{GPS20,GPS24M}。由 \cite{GPS24S} 建立的 \emph{扇形下降}（sectorial descent）应当镜像于凝聚复形的真下降，它可以直接推出环簇退化族中心纤维的同调镜像对称 $\Coh (X_0,J_0)\simeq\mathcal{W} (Y_0,\omega_0)$。Seidel 的提法 \cite{Seidel02} 则提示 $(X_t,J_t)$ 应当镜像于 $Y_0$ 的一个紧化 $(Y,\omega_{p(t)})$，而 $\omega_{p(t)}$ 记录紧化数据附近的体积。

\begin{conceptbox}{扇形下降与真下降}
\emph{扇形下降}（sectorial descent）：一个 Liouville 流形若被一族 Liouville 扇形覆盖，则其包裹 Fukaya 范畴可由各扇形上的范畴沿重叠部分作同伦余极限（粘合）而重建。\emph{真下降}：对代数簇的开覆盖，凝聚层复形由限制到各开集再沿拉回—推出粘合得到，「真」（proper）保证导出前推保持凝聚性。二者互为镜像，是本文「把 A 侧拼装问题化成 B 侧粘合问题」的技术支点。
\par\smallskip\noindent{\footnotesize 参考：nLab，\emph{sectorial descent}；Ganatra--Pardon--Shende \cite{GPS24S}。}
\end{conceptbox}

在沿此方向推进之前，有一个需先检查的合理性问题。给定环簇退化族 $(X_t,J_t)_{t\in\mathbb{C}}$，我们总可以沿具有非平凡对数结构的除子 $Z$ 把中心纤维 $X_0$ 换成它的 \emph{对数爆破} $\Bl_Z (X_0)$。这一操作只与 $0$ 处的纤维有关，不应当改变其光滑化，因此我们自然要问：

\begin{ques}\label{Q:log-blow-up-mirror}
    对数爆破的镜像空间是什么？对数爆破的镜像是否影响紧化？
\end{ques}

只要相信问题 \ref{Q:HMS-toric-degeneration}，问题 \ref{Q:log-blow-up-mirror} 的第二部分就应当是清楚的：它不应当影响紧化。但要回答第二部分，必须先理解第一部分，即构造对数爆破的镜像空间。最简单的情形是，设 $X_0=Z_1\cup_D Z_2$ 为沿除子 $D$ 粘合的两个光滑代数簇，则对数爆破 $\Bl_Z X_0=Z_1\cup\mathbb{P}\left(N_{Z_1} D\oplus N_{Z_2} D\right)\cup Z_2$ 是把 $D$ 换成 $D$ 关于 $Z_1$、$Z_2$ 的法丛直和的射影化得到的。于是只需理解射影丛 $\mathbb{P}\left(N_{Z_1} D\oplus N_{Z_2} D\right)$ 的镜像，再对两个模型应用扇形下降与真下降即可。因此我们真正需要处理的问题是：

\begin{ques}\label{Q:mirror-projective-bundle}
    给定光滑射影代数簇 $X$，对 $X$ 上某条线丛 $\mathcal{L}$，射影化丛 $\mathbb{P} (\mathcal{O}_X\oplus\mathcal{L})$ 的镜像空间是什么？
\end{ques}

\begin{yizhu}[对数爆破为什么会出现在这里]
对数结构记录了 $X_0$ 的奇点以何种方式「展平」成除子；对数爆破则把非正规交叉处沿法方向换成射影丛，是从中心纤维 $X_0$ 走向光滑化 $X_t$ 的局部模型。既然光滑化不依赖中心纤维上的对数结构，辛侧的镜像也应当在相差一个「局部操作」的意义下不变——这正是此处要验证的猜测。注意 $X_0$ 与 $\Bl_Z X_0$ 有相同的光滑化，故二者理应有相同的镜像紧化。
\end{yizhu}

这一问题在某种意义上是局部的。\emph{局部镜像对称} 之名源于 Chiang--Klemm--Yau--Zaslow \cite{CKYZ99} 的工作，原指对 Calabi--Yau 曲面上反典范线丛的全空间作 Gromov--Witten 曲线计数。注意，对概形 $X$ 上任意线丛 $\mathcal{L}_1,\mathcal{L}_2$ 都有 $\mathbb{P}_X (\mathcal{L}_1\oplus\mathcal{L}_2)=\mathbb{P}_X (\mathcal{L}_1\otimes\mathcal{L}_2^{-1}\oplus\mathcal{O}_X)$，因此 \cite{CKYZ99} 所考虑的空间也与本文关心的情况相近。

\begin{conceptbox}{局部镜像对称}
经典局部镜像对称（Chiang--Klemm--Yau--Zaslow）研究 Calabi--Yau 曲面 $S$ 上反典范线丛 $K_S$ 的全空间：一方面它是非紧的 Calabi--Yau 三维流形，其 Gromov--Witten 不变量可用曲线计数定义；另一方面它的镜像是一个复三维的局部模型，其周期积分满足一组 Picard--Fuchs 方程。「局部」指的是只看一个紧化因子的邻域（全空间 $K_S$ 是 $S$ 的管状邻域），与整体 Calabi--Yau 的镜像对称相对。本文把这一计数几何层面的对应提升为范畴层面的同调镜像对称。
\par\smallskip\noindent{\footnotesize 参考：nLab，\emph{local mirror symmetry}；Chiang--Klemm--Yau--Zaslow, \emph{Local mirror symmetry: calculations and interpretations}, 1999。}
\end{conceptbox}

% [原文 p.4]
在陈述主要结果之前，先给出关于 $\mathbb{P} (\mathcal{O}_X\oplus\mathcal{L})$ 的镜像应当是什么样子的直觉。设 $F\to P\to B$ 为一般的全纯 Kähler 丛（纤维与底空间都是 Kähler 流形，并在全空间上给出整体 Kähler 形式），且纤维 $F$ 与底空间 $B$ 分别具有到光滑流形 $\Sigma (F)$、$\Sigma (B)$ 的 Lagrangian 环面纤维化，从而在 $P$ 上诱导整体的环面纤维化。从 SYZ 的角度，镜像应当由对偶 Lagrangian 环面纤维化得到，于是得到纤维丛
\[
    \check{B}\to\check{P}\to\check{F}
\]
其箭头方向反转。因此，我们期望 $\mathbb{P} (\mathcal{O}_X\oplus\mathcal{L})$ 的镜像 $\mathcal{Y}$ 是 $\mathbb{P}^1$ 的镜像上的辛纤维丛；由 \cite{FLTZ12} 知 $\mathbb{P}^1$ 的镜像是 Weinstein 扇形 $(\mathbb{C}^{\ast}, z+1/z)$，其纤维为 $X$ 的镜像 $Y$。给定单值数据 $\phi$ 的 $\mathbb{C}^{\ast}$ 上的辛纤维丛此前已由 Kartal \cite{Kartal21} 研究，他给出了一个漂亮的公式，由 Fukaya 范畴 $\mathcal{W} (Y)$ 及诱导自等价 $\phi\colon\mathcal{W} (Y)\to\mathcal{W} (Y)$ 得到 $\mathcal{W} (\mathcal{Y})$ 的 Fukaya 范畴。因此我们期望：一旦 $\phi$ 镜像于张量乘以 $\mathcal{L}$，带停止位的「映射柱范畴」就拟等价于射影丛 $\mathbb{P}\left(\mathcal{O}_X\oplus\mathcal{L}\right)$ 上的凝聚层复形。不过，这与 \cite{Kartal21} 中提到的镜像对称期望不同。

\begin{conceptbox}{Weinstein 扇形与包裹 Fukaya 范畴}
Liouville 流形指带辛形式 $\omega=d\lambda$、且 $(\omega,\lambda)$ 具有有限型完成化的辛流形；Weinstein 流形是带 Morse 函数的 Liouville 流形，其水平集为接触型。\emph{Weinstein 扇形}是这类流形的带角版本：边界由若干光滑面与角棱组成，可以容许「停止位」（stop）——边界上一族余等向子流形。\emph{包裹 Fukaya 范畴} $\mathcal{W}(X)$ 以紧的 exact Lagrangian 子流形（连同局部系统）为对象，态射由包裹 Floer 上同调给出：其定义中让 Reeb 流沿辛化方向跑到无穷远，停止位的存在则把逃逸路径挡住，得到「部分包裹」版本。$\mathbb{P}^1$ 的镜像 $(\mathbb{C}^{\ast},z+1/z)$ 可看作 $T^{\ast}S^1$ 的扇形模型，其中的超势 $z+1/z$ 记两端各一个停止位。
\par\smallskip\noindent{\footnotesize 参考：nLab，\emph{Weinstein manifold}、\emph{wrapped Fukaya category}；Ganatra--Pardon--Shende \cite{GPS20} 第 2 节。}
\end{conceptbox}

\subsection{主要结果}
本文的主要结果完整回答了问题 \ref{Q:mirror-projective-bundle}，并给出若干推论。设 $Y$ 为 $\mathbb{C}$ 上光滑拟射影代数簇，$L(Y,\xi)^{\times}\to Y$ 为纤维是代数环面 $\mathbb{G}_m$、特征类为 $\xi\in CH^2 (Y)$ 的主丛。设 $L(Y,\xi)\to Y$ 与 $P(Y,\xi)\to Y$ 为相伴的线丛与射影线丛。设 $X$ 为 $Y$ 的镜像辛流形，$\phi\colon X\to X$ 为辛同胚，使得 $\phi\colon\mathcal{W} (X)\to\mathcal{W} (X)$ 镜像于张量乘以线丛 $L(Y,\xi)^{-1}$ 所得的函子 $-\otimes L(Y,\xi)^{-1}$。$\phi$ 的映射柱 $M(X,\phi)\to\mathbb{C}^{\ast}$ 是 $\mathbb{C}^{\ast}$ 上的纤维丛；记 $z$ 与 $1/z$ 为复合函数 $M(X,\phi)\to\mathbb{C}^{\ast}\to\mathbb{C}$。

\begin{yizhu}[映射柱 $M(X,\phi)$ 是什么]
取 $X\times\mathbb{R}$，令 $\mathbb{Z}$ 以 $(x,t)\mapsto(\phi(x),t-1)$ 作用，商空间即 $M(X,\phi)$；它到 $\mathbb{R}/\mathbb{Z}\cong S^1$ 的投影给出纤维丛，纤维为 $X$，绕 $S^1$ 一周的单值性恰为 $\phi$。把 $S^1$ 换成 $\mathbb{C}^{\ast}$（其同伦型相同）并补上复结构，即得文中的 $M(X,\phi)\to\mathbb{C}^{\ast}$；$z$ 是 $\mathbb{C}^{\ast}$ 的坐标函数。当 $\phi=\mathrm{id}$ 时 $M(X,\phi)=X\times\mathbb{C}^{\ast}$，回到平凡的「直积」情形。
\end{yizhu}

\begin{thm}\label{thm:HMS-projective-bundle}
    对 HMS 对 $(X,Y)$ 与镜像自同构 $(\phi,\xi)$（其中 $X$ 为 Weinstein 扇形、$Y$ 光滑拟射影），存在典范的 $A_{\infty}$-函子
\[
        \Coh\left(P(Y,\xi)\right)_{dg}\xrightarrow{\simeq}\Perf\mathcal{W}\left(M(X,\phi),z+\dfrac{1}{z}\right)^{op}
    \]
    它在取同伦范畴后诱导三角范畴的等价。
\end{thm}

\begin{yizhu}[定理 \ref{thm:HMS-projective-bundle} 在说什么]
B 侧是射影丛 $P(Y,\xi)=\mathbb{P}(\mathcal{O}_Y\oplus L(Y,\xi))$ 上凝聚层复形的 dg 范畴；A 侧是映射柱 $M(X,\phi)$ 上带停止位 $z+1/z$ 的包裹 Fukaya 范畴的完美模，取 $op$ 是为了让复合方向与 B 侧一致。几何上：$P^1$ 纤维的镜像正是 $\mathbb{C}^{\ast}$ 方向，$Y$ 的镜像为 $X$，而线丛 $L(Y,\xi)$ 的「扭转」在 A 侧表现为绕 $\mathbb{C}^{\ast}$ 的单值性 $\phi$。Kartal 的公式给出 $\mathcal{W}(M(X,\phi))$ 的生成元与关系，本文进一步证明它与 $\Coh(P(Y,\xi))$ 的 dg 范畴在 $A_{\infty}$ 层面相等。
\end{yizhu}

可以预期，$Y$ 作为零截面或无穷远截面嵌入 $P(Y,\xi)$ 的包含映射 $Y\xrightarrow[i_0]{i_{\infty}} P(Y,\xi)$ 镜像于 Orlov 函子 $\mathcal{W} (X)\xrightarrow[\cup_0]{\cup_{\infty}}\mathcal{W} (M(X,\phi),z+1/z)$，后者分别把 $X$ 送到 $0$ 或 $\infty$ 附近的纤维。在此预期下可以很快得到：

\begin{crl}\label{crl:mirror-symmetry-line-principal-bundles}
    在与定理 \ref{thm:HMS-projective-bundle} 相同的约定下，$A_{\infty}$-范畴之间存在拟等价
\[
        \Coh (L(Y,\xi))_{dg}\xrightarrow{\simeq}\Perf\mathcal{W}\left(M(X,\phi),z\right)^{op}
    \]
    与
\[
        \Coh (L(Y,\xi)^{\times})_{dg}\xrightarrow{\simeq}\Perf\mathcal{W} (M(X,\phi))^{op}
        \]
    它们连接 $Y$ 上 $\mathbb{G}_m$-主丛与线丛的凝聚复形，与建立在 $(X,\phi)$ 之上、可能带一个停止位的映射柱的包裹 Fukaya 范畴上的完美模。
\end{crl}

我们想指出，尽管定理相当一般，我们并不知道很多例子能同时知道同调镜像对称 $\Coh (Y)_{dg}\simeq\Perf\mathcal{W} (X)^{op}$ 以及（分次）辛同胚 $\phi\in\Symp (X)$ 与导出自等价 $\check{\phi}\colon\Coh (Y)\to\Coh (Y)$ 之间的镜像。已知的例子包括环簇 Fano 流形 \cite{FLTZ12,Abouzaid09,Hanlon19,GPS24M} 并张量环簇向量丛，以及对数 Calabi--Yau 曲面 \cite{HK22,HK25} 并张量任意线丛。最近 Varolgunes \cite{Varolgunes26} 进一步推广了 Hacking--Keating 中的同调镜像对称结果，但尚不知道是否有关于镜像自等价群的类似陈述。

% [原文 p.5]
Mamaev 最近的工作 \cite{Mamaev26} 把范畴 $\Coh (P(Y,\xi))_{dg}$ 等同于带 dg 系数的 Fukaya 范畴 $\Perf\mathcal{W} ((\mathbb{C}^{\ast},z+z^{-1});\Coh (Y)_{dg})^{op}$。因此把两个结果结合起来，想必得到拟等价
\[
    \Perf\mathcal{W} (M(X,\phi);z+z^{-1})^{op}\simeq\Perf\mathcal{W} (\mathbb{C}^{\ast},z+z^{-1};\Perf\mathcal{W} (X)^{op})^{op}
\]
这在特殊情形下验证了：$(\mathbb{C}^{\ast}, z+z^{-1})$ 上纤维丛的 Fukaya 范畴应当对应于 Fukaya 范畴的反常 Schober（perverse Schober）。

\begin{yizhu}[「反常 Schober」是什么]
Schober 是把一族范畴沿底空间的粘合数据（「范畴值层」）公理化：底空间每个开集配一个范畴，粘合由拉回—推出给出，单值性则表现为绕圈自等价。B 侧把 $P(Y,\xi)\to Y$ 的 $\mathbb{P}^1$ 纤维方向「摊开」成 $\mathbb{C}^{\ast}$，得到 $\Coh(Y)$ 沿 $\mathbb{C}^{\ast}$ 的范畴族；Mamaev 的定理说这恰好是 $\mathcal{W}((\mathbb{C}^{\ast},z+z^{-1});\Coh(Y))$。本文的主定理则从 A 侧给出同一现象的镜像表述，故两者合起来即上述拟等价。
\end{yizhu}

\begin{keytakeaways}
  \item 出发点：把通常的 $\mathcal{F}(X)\simeq\Coh(Y)$ 反过来看。退化族 $X_t$ 属 B 侧（复几何），镜像 $Y$ 属 A 侧（辛几何），于是问题变成 $\Coh(X_t)$ 是否等价于 $\mathcal{F}(Y)$。
  \item 障碍：对数爆破只改变中心纤维而不改变光滑化，故其镜像应当存在且不影响紧化；最简情形归结为射影丛 $\mathbb{P}(\mathcal{O}_X\oplus\mathcal{L})$ 的镜像问题（问题 \ref{Q:mirror-projective-bundle}）。
  \item 直觉答案：$\mathbb{P}^1$ 的镜像是 Weinstein 扇形 $(\mathbb{C}^{\ast},z+1/z)$，故 $\mathbb{P}(\mathcal{O}\oplus\mathcal{L})$ 的镜像应是 $\mathbb{C}^{\ast}$ 上的辛纤维丛，纤维为 $Y$ 的镜像 $X$，单值性 $\phi$ 镜像于张量乘以 $\mathcal{L}$；Kartal 的公式把它的 Fukaya 范畴归约为 $(X,\phi)$ 的映射柱范畴。
  \item 主定理 \ref{thm:HMS-projective-bundle} 把这套直觉实现为 $A_{\infty}$-等价；推论 \ref{crl:mirror-symmetry-line-principal-bundles} 用 Orlov 函子把零截面/无穷远截面对应到 $0$/$\infty$ 附近的纤维。
\end{keytakeaways}

\begin{yizhu}[两种 Fukaya 范畴记号]
文中 $\mathcal{F}(X,\omega)$ 泛指「某一版本」的 Fukaya 范畴（在引言里保留模糊性：可能是紧致 Fukaya 范畴、可能是精确 Fukaya 范畴）；而 $\mathcal{W}(X)$ 专指包裹 Fukaya 范畴。本文 A 侧的对象几乎都是 Weinstein 扇形，所以主定理与推论中一律用 $\mathcal{W}$；只有在回顾背景或讨论一般镜像对时才用 $\mathcal{F}$。
\end{yizhu}
